\chapter{Especificação da ALU}
\label{cap:especificacao}

A ALU realiza operações lógicas e aritméticas em palavras de 8 bits e executa 16
funções diferentes, entre operações lógicas, aritméticas, de rotação e de
deslocamento. O circuito é totalmente combinacional \cite{tp1}. A
Figura~\ref{fig:bloco} mostra o diagrama de blocos com as entradas e saídas.

\begin{figure}[htb]
\caption{Diagrama de blocos da ALU.}\label{fig:bloco}
\centering
\begin{tikzpicture}[
  >={Stealth[length=2mm]},
  port/.style={font=\small\ttfamily}]
  \node[draw, minimum width=3cm, minimum height=4cm, font=\large\bfseries] (alu) {ALU};
  % entradas
  \draw[->, line width=1.2pt] ($(alu.west)+(-2.2,1.4)$) -- ($(alu.west)+(0,1.4)$);
  \node[port, anchor=east] at ($(alu.west)+(-2.3,1.4)$) {a\_in[7..0]};
  \draw[->, line width=1.2pt] ($(alu.west)+(-2.2,0.5)$) -- ($(alu.west)+(0,0.5)$);
  \node[port, anchor=east] at ($(alu.west)+(-2.3,0.5)$) {b\_in[7..0]};
  \draw[->] ($(alu.west)+(-2.2,-0.5)$) -- ($(alu.west)+(0,-0.5)$);
  \node[port, anchor=east] at ($(alu.west)+(-2.3,-0.5)$) {c\_in};
  \draw[->, line width=1.2pt] ($(alu.west)+(-2.2,-1.4)$) -- ($(alu.west)+(0,-1.4)$);
  \node[port, anchor=east] at ($(alu.west)+(-2.3,-1.4)$) {op\_sel[3..0]};
  % saídas
  \draw[->, line width=1.2pt] ($(alu.east)+(0,0.8)$) -- ($(alu.east)+(2.2,0.8)$);
  \node[port, anchor=west] at ($(alu.east)+(2.3,0.8)$) {r\_out[7..0]};
  \draw[->, line width=1.2pt] ($(alu.east)+(0,-0.8)$) -- ($(alu.east)+(2.2,-0.8)$);
  \node[port, anchor=west] at ($(alu.east)+(2.3,-0.8)$) {flags\_out[4..0]};
\end{tikzpicture}
\fonte{Elaborada pelo autor, com base em \textcite{tp1}.}
\end{figure}

\section{Entrada e Saida}

\begin{quadro}[htb]
  \caption{Portas da ALU.}
  \label{quad:portas}
  \centering
  \begin{tabular}{lcl}
    \hline
    Porta & Direção & Descrição \\ \hline
    \texttt{a\_in[7..0]} & entrada & Operando A \\
    \texttt{b\_in[7..0]} & entrada & Operando B (operações com dois operandos) \\
    \texttt{c\_in} & entrada & Carry de entrada (somas, subtrações com carry e rotações) \\
    \texttt{op\_sel[3..0]} & entrada & Seleção da operação \\
    \texttt{r\_out[7..0]} & saída & Resultado \\
    \texttt{flags\_out[4..0]} & saída & Flags: [0] N, [1] V, [2] Z, [3] H, [4] C \\ \hline
  \end{tabular}
  \fonte{Especificação do primeiro projeto (TP1).}
\end{quadro}

\begin{quadro}[htb]
\caption{\emph{Flags} de saída da ALU.}\label{qua:flags}
\small
\begin{tabular}{|l|l|p{8.2cm}|}
\hline
\textbf{Bit} & \textbf{Nome} & \textbf{Significado} \\ \hline
\texttt{flags\_out[0]} & \emph{negative} (\texttt{n\_flag}) & Bit 7 do resultado. Nas operações aritméticas indica resultado positivo (\texttt{'0'}) ou negativo (\texttt{'1'}). \\ \hline
\texttt{flags\_out[1]} & \emph{overflow} (\texttt{v\_flag}) & Vale \texttt{'1'} quando há \emph{overflow} em soma ou subtração. \\ \hline
\texttt{flags\_out[2]} & \emph{zero} (\texttt{z\_flag}) & Vale \texttt{'1'} quando o resultado é zero. \\ \hline
\texttt{flags\_out[3]} & \emph{half carry/borrow} (\texttt{h\_flag}) & Na soma, é o \emph{carry out} do bit 3 para o bit 4; na subtração, é o \emph{borrow out}. \\ \hline
\texttt{flags\_out[4]} & \emph{carry/borrow} (\texttt{c\_flag}) & Na soma, é o \emph{carry out} do bit mais significativo; na subtração, é o \emph{borrow out}. Também é usada nas operações de rotação. \\ \hline
\end{tabular}
\fonte{Adaptado de \textcite{tp1}.}
\end{quadro}

\newpage
\section{Operações}

\begin{longtable}{c l p{7.2cm} c}
\caption{Operações da ALU.}\label{tab:operacoes}\\
\hline \texttt{op\_sel} & Mnemônico & Operação & Flags C / H / V \\ \hline \endfirsthead
\hline \texttt{op\_sel} & Mnemônico & Operação & Flags C / H / V \\ \hline \endhead
0000 & PASS\_B & $r = b$ & 0 / 0 / -- \\
0001 & OR & $r = a \lor b$ & 0 / 0 / -- \\
0010 & XOR & $r = a \oplus b$ & 0 / 0 / -- \\
0011 & AND & $r = a \land b$ & 0 / 0 / -- \\
0100 & ADDC & $r = a + b + c_{in}$ & carry / half carry / ovf \\
0101 & ADD & $r = a + b$ & carry / half carry / ovf \\
0110 & SUBC & $r = a - b - c_{in}$ & borrow / half borrow / ovf \\
0111 & SUB & $r = a - b$ & borrow / half borrow / ovf \\
1000 & NOT & $r = \overline{a}$ & 0 / -- / -- \\
1001 & SRA & $r = a_7, a_{7..1}$ (desloc. aritmético à direita) & $a_0$ / -- / -- \\
1010 & SRL & $r = 0, a_{7..1}$ (desloc. lógico à direita) & $a_0$ / -- / -- \\
1011 & SLL & $r = a_{6..0}, 0$ (desloc. lógico à esquerda) & $a_7$ / -- / -- \\
1100 & RRC & $r = c_{in}, a_{7..1}$ (rotação à direita pelo carry) & $a_0$ / -- / -- \\
1101 & RLC & $r = a_{6..0}, c_{in}$ (rotação à esquerda pelo carry) & $a_7$ / -- / -- \\
1110 & RR & $r = a_0, a_{7..1}$ (rotação à direita) & $a_0$ / -- / -- \\
1111 & RL & $r = a_{6..0}, a_7$ (rotação à esquerda) & $a_7$ / -- / -- \\ \hline
\multicolumn{4}{p{13cm}}{\footnotesize Em todas as operações, Z = `1' se $r = 0$ e N = $r_7$. O símbolo ``--'' indica \textit{don't care}.} \\
\end{longtable}

As operações de deslocamento e rotação são ilustradas na
Figura~\ref{fig:shift}, em que \textbf{C} é o \emph{carry}.

\begin{figure}[htb]
\caption{Operações de deslocamento e rotação.}\label{fig:shift}
\centering
\begin{tikzpicture}[
  >={Stealth[length=2mm]},
  reg/.style={draw, minimum width=2.6cm, minimum height=6mm, font=\small},
  bit/.style={draw, minimum width=6mm, minimum height=6mm, font=\small},
  lab/.style={font=\small\bfseries, anchor=east}]
  % SRA
  \node[lab] at (-3.2,0) {SRA};
  \node[reg] (r0) at (0,0) {b7 \dots\ b0};
  \node[bit] (c0) at (2.2,0) {C};
  \draw[->] (r0.east) -- (c0.west);
  \draw[->] ($(r0.north west)+(3mm,0)$) -- ++(0,4mm) -- ++(-6mm,0) |- (r0.west);
  % SRL
  \node[lab] at (-3.2,-1.1) {SRL};
  \node[bit] (z1) at (-2,-1.1) {0};
  \node[reg] (r1) at (0,-1.1) {b7 \dots\ b0};
  \node[bit] (c1) at (2.2,-1.1) {C};
  \draw[->] (z1) -- (r1.west); \draw[->] (r1.east) -- (c1.west);
  % SLL
  \node[lab] at (-3.2,-2.2) {SLL};
  \node[bit] (c2) at (-2,-2.2) {C};
  \node[reg] (r2) at (0,-2.2) {b7 \dots\ b0};
  \node[bit] (z2) at (2.2,-2.2) {0};
  \draw[<-] (c2.east) -- (r2.west); \draw[<-] (r2.east) -- (z2.west);
  % RRC
  \node[lab] at (-3.2,-3.7) {RRC};
  \node[reg] (r3) at (0,-3.7) {b7 \dots\ b0};
  \node[bit] (c3) at (2.2,-3.7) {C};
  \draw[->] (r3.east) -- (c3.west);
  \draw[->] (c3.south) -- ++(0,-4mm) -| ($(r3.west)+(-5mm,0)$) -- (r3.west);
  % RLC
  \node[lab] at (-3.2,-5.2) {RLC};
  \node[bit] (c4) at (-2,-5.2) {C};
  \node[reg] (r4) at (0,-5.2) {b7 \dots\ b0};
  \draw[<-] (c4.east) -- (r4.west);
  \draw[->] (c4.south) -- ++(0,-4mm) -| ($(r4.east)+(5mm,0)$) -- (r4.east);
  % RR
  \node[lab] at (-3.2,-6.7) {RR};
  \node[reg] (r5) at (0,-6.7) {b7 \dots\ b0};
  \node[bit] (c5) at (2.2,-6.7) {C};
  \draw[->] (r5.east) -- (c5.west);
  \draw[->] ($(r5.east)+(3mm,0)$) -- ++(0,-5mm) -| ($(r5.west)+(-5mm,0)$) -- (r5.west);\fill ($(r5.east)+(3mm,0)$) circle (0.6pt);
  % RL
  \node[lab] at (-3.2,-8.2) {RL};
  \node[bit] (c6) at (-2,-8.2) {C};
  \node[reg] (r6) at (0,-8.2) {b7 \dots\ b0};
  \draw[<-] (c6.east) -- (r6.west);
  \draw[->] ($(r6.west)+(-2mm,0)$) -- ++(0,-5mm) -| ($(r6.east)+(5mm,0)$) -- (r6.east);\fill ($(r6.west)+(-2mm,0)$) circle (0.6pt);
\end{tikzpicture}
\fonte{Adaptado de \textcite{tp1}.}
\end{figure}
